\documentclass[aspectratio=169,11pt]{beamer}
\usetheme{Madrid}
\usecolortheme{default}
\definecolor{navy}{HTML}{142D4E}
\definecolor{teal}{HTML}{087F8C}
\definecolor{pale}{HTML}{EAF6F6}
\setbeamercolor{structure}{fg=teal}
\setbeamercolor{frametitle}{fg=white,bg=navy}
\setbeamercolor{title}{fg=white}
\setbeamercolor{titlelike}{fg=white,bg=navy}
\setbeamercolor{block title}{fg=white,bg=teal}
\setbeamercolor{block body}{fg=black,bg=pale}
\setbeamertemplate{navigation symbols}{}
\setbeamertemplate{footline}[frame number]
\usepackage[T1]{fontenc}
\usepackage{lmodern}
\usepackage{listings}
\usepackage{xcolor}
\usepackage{colortbl}
\lstset{basicstyle=\ttfamily\small,backgroundcolor=\color{pale},frame=single,rulecolor=\color{teal},columns=fullflexible,keepspaces=true,showstringspaces=false}
\title{Applications of Bitwise Operators}
\subtitle{IPv4 Addresses, Subnet Masks, and Networks}
\author{COMP 211 --- Systems Programming}
\date{}
\begin{document}

\begin{frame}
\titlepage
\begin{center}\large How does a computer figure out which network an IP address belongs to?\end{center}
\end{frame}

\begin{frame}{Networks, Hosts, and IP Addresses}
\begin{block}{Think of a neighborhood}
\begin{itemize}
\item \textbf{Network:} a group of devices that can communicate.
\item \textbf{Host:} a device or interface on that network.
\item \textbf{IP address:} a numerical address used to deliver data.
\end{itemize}
\end{block}
\vspace{0.4cm}
\begin{center}
\Large \texttt{192.168.1.10}\par\medskip
\normalsize Neighborhood = network \quad $\vert$ \quad House = host
\end{center}
\end{frame}

\begin{frame}{IPv4 and Subnet Masks}
\begin{itemize}
\item An IPv4 address contains \textbf{32 bits}: four groups of 8 bits.
\item The \textbf{subnet mask} says which bits identify the network.
\item In \texttt{192.168.1.10/24}, the first 24 bits are network bits.
\end{itemize}
\vspace{0.3cm}
\begin{center}
\begin{tabular}{|c|c|c|c|}\hline
\cellcolor{pale}192 & \cellcolor{pale}168 & \cellcolor{pale}1 & 10\\\hline
\multicolumn{3}{|c|}{\textbf{24 network bits}} & \textbf{8 host bits}\\\hline
\end{tabular}
\end{center}
\vspace{0.35cm}
\begin{block}{Subnet mask}\centering\texttt{255.255.255.0} \quad = \quad \texttt{/24}\end{block}
\end{frame}

\begin{frame}[fragile]{Bitwise AND Finds the Network}
\begin{block}{Rule}
\centering \textbf{Network address = IP address \texttt{\&} subnet mask}
\end{block}
\begin{lstlisting}
IP:       192.168.1.10
Mask:     255.255.255.0
Result:   192.168.1.0

Last octet:
  00001010   (10)
& 00000000   (0)
----------
  00000000   (0)
\end{lstlisting}
\centering \textbf{AND keeps network bits and clears host bits.}
\end{frame}

\begin{frame}{Subnetting: Splitting a Network}
\begin{itemize}
\item Start with \texttt{192.168.1.0/24}: \textbf{256 total addresses}.
\item Borrow \textbf{one host bit} to change \texttt{/24} into \texttt{/25}.
\item Now we have \textbf{two subnets}, each with 128 addresses.
\end{itemize}
\vspace{0.35cm}
\begin{block}{Two /25 subnets}
\textbf{Subnet A:} \texttt{192.168.1.0 -- 192.168.1.127}\\[6pt]
\textbf{Subnet B:} \texttt{192.168.1.128 -- 192.168.1.255}
\end{block}
\vspace{0.25cm}
\centering Last octet starts with \texttt{0xxxxxxx} or \texttt{1xxxxxxx}.
\end{frame}

\begin{frame}[fragile]{Broadcast Addresses: All Host Bits Are 1}
\begin{itemize}
\item \textbf{Network address:} all host bits are \texttt{0}.
\item \textbf{Broadcast address:} all host bits are \texttt{1}.
\end{itemize}
\begin{lstlisting}
Network:         192.168.1.0
Last usable:     192.168.1.254
Broadcast:       192.168.1.255
\end{lstlisting}
\begin{block}{Using OR and NOT in C}
\centering \texttt{broadcast = ip | \~{}subnet\_mask;}
\end{block}
\small For ordinary IPv4 subnets, the broadcast address is not assigned to an individual host.
\end{frame}

\begin{frame}{Audience Question}
\Large Are these computers on the same \texttt{/24} subnet?\par\vspace{0.45cm}
\begin{center}
\Large \textbf{A:} \texttt{192.168.1.10}\\[0.3cm]
\Large \textbf{B:} \texttt{192.168.1.50}
\end{center}
\vspace{0.5cm}
\pause
\begin{block}{Answer: Yes!}
Both addresses ANDed with \texttt{255.255.255.0} produce \texttt{192.168.1.0}.
\end{block}
\end{frame}

\begin{frame}{Key Takeaways}
\begin{itemize}
\setlength{\itemsep}{0.55em}
\item An IPv4 address contains \textbf{network and host bits}.
\item A \textbf{subnet mask} tells us how those bits are divided.
\item \textbf{AND} calculates the network address.
\item \textbf{Subnetting} creates smaller networks by borrowing host bits.
\item \textbf{OR and NOT} help calculate the broadcast address.
\end{itemize}
\vspace{0.35cm}
\begin{block}{Connection to COMP 211}
Simple binary operations help computers make real networking decisions.
\end{block}
\end{frame}
\end{document}
